\section{Аннотация}
Разработана учебная система подбора видеоигр на основе логического представления знаний.
Предметная область включает игровые режимы, сюжетность, сложность, независимость разработки,
жанры и студии. База Prolog содержит 51 унарный факт, 15 бинарных фактов и семь правил.
На её основе построена OWL-онтология с ограничениями, выводимыми классами и правилом SWRL.
Консольная система поддержки принятия решений уточняет предпочтения и обращается к SWI-Prolog.
Проверены положительные и отрицательные сценарии, а также соответствие результатов Prolog и HermiT.
Основное ограничение решения связано с неполнотой учебной базы и различием предположений
закрытого и открытого мира.

\section{Введение и анализ требований}
Цель модуля --- освоить логическое моделирование, онтологии и проверяемый вывод,
а затем использовать их в рекомендательном приложении. В отличие от модели,
обученной на пользовательских оценках, здесь результат определяется явными фактами
и правилами. Это позволяет объяснить рекомендацию и воспроизвести её при тех же данных.
Один набор знаний используется в Prolog, в онтологии и в диалоговой программе.

Нумерация работ принята по выданному заданию «Системы ИИ»: лабораторная 1 включает
базу знаний и её онтологическое представление, лабораторная 2 --- DSS.
В пособии ИТМО этим этапам соответствуют работы 1, 2 и 3 на страницах 56--58.
Постановка задачи и структура расширенного отчёта также учитывают приложение 2
пособия и шаблон отчёта из задания.

Требуется не менее 20 унарных фактов, 10--15 бинарных фактов и 5--7 правил;
нужны запросы разной сложности и комментарии. Онтология должна содержать классы,
свойства, индивиды, иерархию, ограничения и демонстрацию вывода.
Для DSS необходимы фиксированный формат входа, короткий диалог, обращение к знаниям,
рекомендации с обоснованием и обработка некорректных данных.

База представляет учебный срез из 12 игр. Признаки «сложная» и «спокойная»
являются принятыми условиями примера, а не результатами объективного измерения.
Программа не оценивает возрастные рейтинги, актуальные цены и совместимость оборудования.
